\scanpage{264}
\mspage{5}

les sommets, suivant que la chaîne est directe ou rétrograde.

Notons que tout \inserted{simplexe type $\unsure{i|j}\geq0$} $\mathbb{\Delta}_i$ peut être vu comme un intervalle, de longueur $i$, de longueur utile $1$ (si $i\geq1$, $0$ si $i=0$). Son origine est $0$, son extrémité est $i$.

Mais appelons une catégorie quelconque une \deleted{intervalle} \uline{catégorie} \uline{intervalle} si elle est isomorphe à une \deleted{catégorie} intervalle standard. Notons ici

\uline{Proposition 1} Soit $I$ un intervalle \unsure{standart|standard} de longueur $n\geq1$. \deleted{Alors origine et extrémité \unsure{coïncident|coïncidant}}

a) \deleted{Considérons} Soit $i\in I$. Pour que \deleted{$i\in\{0,n\}$ i.e.} \inserted{$i\in\{0,n\}$, i.e.} pour que $i$ soit un des deux ``\uline{bouts}'' \inserted{(ou ``\uline{extrémités}'')} (source ou but) de $I$, il f. et s. \unsure{ou|que} l'une des deux catégories $I_{/i}$ ou $i\backslash I$ soit \deleted{\unsure{triviale}} la catégorie ponctuelle \inserted{et que l'autre soit iso. à un $\mathbb{\Delta}_j$ $*$} (Caractérisation intrinsèque des bouts d'un intervalle.) Pour que $i$ soit la source, il f. et s. que \unsure{$i$ est}

% slanted
\margin{$*$ où $j$ est la longueur de l'unique chaîne structurale qui ``touche'' $i$ (i.e. telle que $i$ soit \deleted{dans} \unsure{la|$\in$} sous-catégorie pleine \unsure{''} définie par ladite chaîne)}

b) Tout foncteur $I\longrightarrow I$ qui fixe les deux \unsure{extrémités} de $I$ est l'identité.

\deleted{c) \unsure{Il en}}

\uline{Corollaire} Le groupe des automorphismes de $I$ a au plus deux éléments. S'il en a deux, l'autom. non trivial échange les deux extrémités. (On dit \deleted{\unsure{que}} alors que l'intervalle est \uline{symétrique}.) \deleted{\unsure{on a trivial}}

\deleted{\uline{Proposition 2}. Si $I$ est un intervalle standard \deleted{\unsure{$I$}} il existe un unique couple $(I',\Theta)$, \unsure{avec|où} $I'$ un \deleted{est isomorphe à un intervalle standard unique} intervalle standard, et $\Theta$ un isomorphisme de $I$ sur $I'$, qui transforme but en origine et origine en but)

$I'$ est déduit de $I$, en transformant, par la symétrie}

\[\deleted{i\longmapsto n-i\qquad [0,n]\longrightarrow[0,n]\,,}\]

% archive sheet number: 263
